%\documentclass[stu, 12pt]{apa7}
\documentclass[stu, 12pt, floatsintext]{apa7}
%\usepackage{fontspec}
%\setmainfont{Times New Roman} % Or your system's equivalent serif font
\usepackage[american]{babel}
\usepackage{csquotes}
\usepackage[style=apa, backend=biber]{biblatex}
%\usepackage[style=authoryear, backend=biber]{biblatex} % Changed from style=apa
\usepackage{float}
\usepackage{multicol}
\usepackage{hyperref} % nameref is bundled with hyperref
\usepackage{graphicx}   % Required for including images
\usepackage{subcaption} % Required for subfigures and subcaptions
%\usepackage{listings}
%\usepackage{xcolor}
\usepackage{array}
\usepackage{booktabs}
\usepackage{multirow}
\usepackage[dvipsnames]{xcolor}
\usepackage{amsmath}

% Define the console style before using it
%\lstdefinestyle{console}{
	%	backgroundcolor=\color{black!5},
	%	basicstyle=\ttfamily\small,
	%	breaklines=true,
	%	frame=single,
	%	rulecolor=\color{black!30}
	%}
\graphicspath{
	{../img/}
	%   {../img_ipynb/}
	%	{../img/H1/}
	%	{../../results/DayResult/}
}
\DeclareGraphicsExtensions{.pdf,.jpeg,.png,.jpg}

\hypersetup{
	colorlinks=true,
	linkcolor=blue,
	citecolor=blue,
	urlcolor=blue
}

\addbibresource{../ref/references.bib}
\title{Anime Face Generation with VAE}
\author{Max Raxy}
\affiliation{Udacity Institute of AI \& Technology}
\course{Capstone-5: Generative AI Applications}
\professor{~}
\duedate{\today}

\begin{document}
	\maketitle
	\tableofcontents
	\newpage
	\listoffigures%\thispagestyle{fancy}
	\listoftables%\thispagestyle{fancy}
	\newpage

\section{Overview}
We are using \textbf{Anime Face Dataset} from \textbf{Kaggle}~\parencites{kaggle-animeface}{org-animefacedataset}. The dataset contains $63565$ anime-face images, and the goal is to use Variational Auto-Encoder (VAE),~\parencites{org-vae1}{org-vae2} to generate anime faces.

\section{Dataset or Prompt Description}
Describe the data or prompt source, its size or scope, and how it supports the generative task.
Dataset contains anime faces with aspect-ratio $1$  (with a very slight deviation from 1: mean=$1.000003$, standard deviation  =$0.000171$). The height and weight varies between $25$ to $220$ (Fig.~\ref{fig:dim-dist})

\begin{figure}[!tbh]
	\centering
	\includegraphics[width=1.0\textwidth]{hw-hist}
	\captionsetup{justification=centering}
	\caption{Width and Height Distributions.}
	\label{fig:dim-dist}
\end{figure}

We decided to resize all images to 64x64  for this project, this size still has good quality, and proper for training on home-PC with limited resources. Sample of original images resized to 64x64 are shown in Fig.~\ref{fig:sample-org}

\begin{figure}[!tbh]
	\centering
	\includegraphics[width=1.0\textwidth]{sample-img-org}
	\captionsetup{justification=centering}
	\caption{Sample original images resized to 64x64.}
	\label{fig:sample-org}
\end{figure}


\section{Model Design and Training Approach}
Explain the model architecture, training setup, and key design decisions.

\section{Output Evaluation and Interpretation}
Discuss generated outputs, highlighting quality, variability, and limitations.

\section{Ethical Considerations and Responsible Use}
Identify at least one ethical concern related to bias, misuse, creative ownership, or societal impact.

\section{Limitations and Future Improvements}
Describe limitations of the model or data and propose realistic next steps.

% apa does not include ref and appx in TOC we forced it
\phantomsection % Ensures hyperref links to the correct page
\addcontentsline{toc}{section}{References}
\printbibliography %[heading=none] % Removes default unnumbered header so it doesn't double-up

% Appendices start here (AFTER references)
\appendix
\phantomsection
\addcontentsline{toc}{section}{Appendix A: Evaluation Metrics for Binary Classification}
\section{Appendix A: Evaluation Metrics for Binary Classification}\label{app:a}
\subsection{ELBO}
In VAEs, the overall objective is to minimize the negative Evidence Lower Bound (ELBO):

$$\mathcal{L}_{\text{total}} = \mathcal{L}_{\text{rec}} + \beta \cdot \mathcal{L}_{\text{KLD}}$$

Here, the reconstruction loss measures how accurately the decoder reconstructs the input (using mean squared error),
while the Kullback-Leibler Divergence ($\text{KLD}$) forces
the latent space representation to approximate a standard normal distribution $\mathcal{N}(0, \mathbf{I})$.

%\newpage
%
%\phantomsection
%\addcontentsline{toc}{section}{Appendix B: XXX}
%\section{Appendix B: XXX}
%Content of Appendix B...
\end{document}
